% GENERATED FILE. Edit canonical lessons, then run npm run generate:books.
% canonical-source-hash: fnv1a64:bb9a36c3fe2d0557
% canonical-lessons: HI-C177-sughna, HI-C177-manana, HI-C177-pisna, HI-C177-ubalna, HI-C177-talna

\chapter{To smell, To celebrate, To grind, To boil, To fry}
\label{ch:hi-to-smell-to-celebrate-to-grind-to-boil-to-fry}
\hlchaptermodality{177}

\begin{chapteropening}
\textbf{By the end of this chapter:} \emph{I can say to smell, to celebrate, to grind, to boil and to fry.}

Complete the last lesson of chapter 177: I can say to smell, to celebrate, to grind, to boil and to fry.
\end{chapteropening}

\section[\texorpdfstring{\hi{सूँघना} (sū̃ghnā)}{सूँघना (sū̃ghnā)}]{\hi{सूँघना} (sū̃ghnā) — to smell}

\label{lesson:HI-C177-sughna}



\begin{quote}
Before the new one: say the Hindi for to throw, then the Hindi for to taste.
\end{quote}

\subsection*{You'll want to know: \hi{सूँघना}}
\textbf{\hi{सूँघना}} — \emph{sū̃ghnā} — \textquotedblleft{}to smell\textquotedblright{}.

\begin{cousinweb}
\textbf{\hi{सूँघना}} (\emph{sūṁghnā}) is a tadbhava, a word inherited from Sanskrit and reshaped by sound change. It comes from \emph{śiṅghati}, \textquotedblleft{}smells\textquotedblright{} (CDIAL 12579). The Sanskrit verb is known only from a list of verb roots.
\end{cousinweb}

\begin{grammarlens}[title={What you've built}]
One more word for this chapter.
\end{grammarlens}

\subsection*{Your turn}
\begin{itemize}
\raggedright
  \item \emph{Say it:} \emph{sū̃ghnā}
  \item \emph{Say it:} \emph{sū̃ghnā}, once more
  \item \emph{Say it:} say \emph{sū̃ghnā}
  \item \emph{Recall:} say the Hindi for to throw, then the Hindi for to taste, then say \emph{sū̃ghnā} again
\end{itemize}

\begin{tcolorbox}[breakable,colback=teal!4,colframe=teal!35!black,title={Before you move on}]
What does \hi{सूँघना} mean? (\textquotedblleft{}To smell\textquotedblright{}.) Say it once more.
\end{tcolorbox}
\section[\texorpdfstring{\hi{मनाना} (manānā)}{मनाना (manānā)}]{\hi{मनाना} (manānā) — to celebrate, to persuade}

\label{lesson:HI-C177-manana}



\begin{quote}
Before the new one: say the Hindi for to taste, then the Hindi for to smell.
\end{quote}

\subsection*{You'll want to know: \hi{मनाना}}
\textbf{\hi{मनाना}} — \emph{manānā} — \textquotedblleft{}to celebrate, to persuade\textquotedblright{}.

\begin{cousinweb}
\textbf{\hi{मनाना}} (\emph{manānā}) is the causing form of \emph{mānnā}, \textquotedblleft{}to accept, honour\textquotedblright{}: to persuade is to make someone accept, and to celebrate is to honour a day. CDIAL lists both under Sanskrit \emph{manyate}, \textquotedblleft{}thinks, honours\textquotedblright{} (CDIAL 9857).
\end{cousinweb}

\begin{grammarlens}[title={What you've built}]
One more word for this chapter.
\end{grammarlens}

\subsection*{Your turn}
\begin{itemize}
\raggedright
  \item \emph{Say it:} \emph{manānā}
  \item \emph{Say it:} \emph{manānā}, once more
  \item \emph{Say it:} say \emph{manānā}
  \item \emph{Recall:} say the Hindi for to taste, then the Hindi for to smell, then say \emph{manānā} again
\end{itemize}

\begin{tcolorbox}[breakable,colback=teal!4,colframe=teal!35!black,title={Before you move on}]
What does \hi{मनाना} mean? (\textquotedblleft{}To celebrate, to persuade\textquotedblright{}.) Say it once more.
\end{tcolorbox}
\section[\texorpdfstring{\hi{पीसना} (pīsnā)}{पीसना (pīsnā)}]{\hi{पीसना} (pīsnā) — to grind}

\label{lesson:HI-C177-pisna}



\begin{quote}
Before the new one: say the Hindi for to smell, then the Hindi for to celebrate.
\end{quote}

\subsection*{You'll want to know: \hi{पीसना}}
\textbf{\hi{पीसना}} — \emph{pīsnā} — \textquotedblleft{}to grind\textquotedblright{}.

\begin{cousinweb}
\textbf{\hi{पीसना}} (\emph{pīsnā}) is a tadbhava, a word inherited from Sanskrit and reshaped by sound change. It comes from \emph{piṁṣati}, \textquotedblleft{}grinds\textquotedblright{} (CDIAL 8142).
\end{cousinweb}

\begin{grammarlens}[title={What you've built}]
One more word for this chapter.
\end{grammarlens}

\subsection*{Your turn}
\begin{itemize}
\raggedright
  \item \emph{Say it:} \emph{pīsnā}
  \item \emph{Say it:} \emph{pīsnā}, once more
  \item \emph{Say it:} say \emph{pīsnā}
  \item \emph{Recall:} say the Hindi for to smell, then the Hindi for to celebrate, then say \emph{pīsnā} again
\end{itemize}

\begin{tcolorbox}[breakable,colback=teal!4,colframe=teal!35!black,title={Before you move on}]
What does \hi{पीसना} mean? (\textquotedblleft{}To grind\textquotedblright{}.) Say it once more.
\end{tcolorbox}
\section[\texorpdfstring{\hi{उबालना} (ubālnā)}{उबालना (ubālnā)}]{\hi{उबालना} (ubālnā) — to boil}

\label{lesson:HI-C177-ubalna}



\begin{quote}
Before the new one: say the Hindi for to celebrate, then the Hindi for to grind.
\end{quote}

\subsection*{You'll want to know: \hi{उबालना}}
\textbf{\hi{उबालना}} — \emph{ubālnā} — \textquotedblleft{}to boil\textquotedblright{}.

\begin{cousinweb}
\textbf{\hi{उबालना}} (\emph{ubālnā}) is a tadbhava, a word inherited and reshaped by sound change. It comes from \emph{*ubbāl}, \textquotedblleft{}make boil\textquotedblright{}, a form CDIAL reconstructs from its descendants (CDIAL 2339). Morgenstierne traced it to \emph{ud-}, \textquotedblleft{}up\textquotedblright{}, and an Indo-European root \textquotedblleft{}to boil\textquotedblright{}; CDIAL finds the forms hard to fit and notes other possible sources.
\end{cousinweb}

\begin{grammarlens}[title={What you've built}]
One more word for this chapter.
\end{grammarlens}

\subsection*{Your turn}
\begin{itemize}
\raggedright
  \item \emph{Say it:} \emph{ubālnā}
  \item \emph{Say it:} \emph{ubālnā}, once more
  \item \emph{Say it:} say \emph{ubālnā}
  \item \emph{Recall:} say the Hindi for to celebrate, then the Hindi for to grind, then say \emph{ubālnā} again
\end{itemize}

\begin{tcolorbox}[breakable,colback=teal!4,colframe=teal!35!black,title={Before you move on}]
What does \hi{उबालना} mean? (\textquotedblleft{}To boil\textquotedblright{}.) Say it once more.
\end{tcolorbox}
\section[\texorpdfstring{\hi{तलना} (talnā)}{तलना (talnā)}]{\hi{तलना} (talnā) — to fry}

\label{lesson:HI-C177-talna}



\begin{quote}
Before the new one: say the Hindi for to grind, then the Hindi for to boil.
\end{quote}

\subsection*{You'll want to know: \hi{तलना}}
\textbf{\hi{तलना}} — \emph{talnā} — \textquotedblleft{}to fry\textquotedblright{}.

\begin{cousinweb}
\textbf{\hi{तलना}} (\emph{talnā}) is a tadbhava, a word inherited and reshaped by sound change. It comes from \emph{*talati}, \textquotedblleft{}fries\textquotedblright{}, a form CDIAL reconstructs from its descendants (CDIAL 5736).
\end{cousinweb}

\begin{grammarlens}[title={What you've built}]
One more word for this chapter.
\end{grammarlens}

\subsection*{Your turn}
\begin{itemize}
\raggedright
  \item \emph{Say it:} \emph{talnā}
  \item \emph{Say it:} \emph{talnā}, once more
  \item \emph{Say it:} say \emph{talnā}
  \item \emph{Recall:} say the Hindi for to grind, then the Hindi for to boil, then say \emph{talnā} again
\end{itemize}

\begin{tcolorbox}[breakable,colback=teal!4,colframe=teal!35!black,title={Before you move on}]
What does \hi{तलना} mean? (\textquotedblleft{}To fry\textquotedblright{}.) Say it once more.
\end{tcolorbox}
